\documentclass[12pt,a4paper]{ctexart}
\usepackage[utf8]{inputenc}
\usepackage{graphicx}
\usepackage{amsmath,amssymb,amsfonts}
\usepackage{geometry}
\usepackage{hyperref}
\usepackage{caption}
\usepackage{subcaption}
\usepackage{enumerate}
\usepackage{algorithm}
\usepackage{algorithmicx}
\usepackage{algpseudocode}
\usepackage{natbib}
\usepackage{xcolor}
\usepackage{listings}

% 页面设置
\geometry{margin=1in}

% 定理环境定义
\newtheorem{theorem}{定理}[section]
\newtheorem{lemma}[theorem]{引理}
\newtheorem{corollary}[theorem]{推论}
\newtheorem{definition}[theorem]{定义}

% 代码样式设置
\lstset{
    language=Python,
    basicstyle=\ttfamily\small,
    keywordstyle=\color{blue},
    stringstyle=\color{red},
    commentstyle=\color{green!60!black},
    numbers=left,
    numberstyle=\tiny\color{gray},
    frame=single,
    breaklines=true
}

% 标题信息
\title{全面的LaTeX语法示例文档}
\author{
    张三 \\
    \small 计算机科学系，XYZ大学 \\
    \small \href{mailto:zhangsan@example.com}{zhangsan@example.com}
    \and
    李四 \\
    \small 数学研究所，ABC研究中心 \\
    \small \href{mailto:lisi@example.com}{lisi@example.com}
}
\date{\today}
\begin{document}

\maketitle
\begin{abstract}
本文档展示了各种核心LaTeX语法元素，包括文档结构、数学公式、图形插入、列表环境、算法伪代码和参考文献引用。通过这个全面的例子，可以快速了解LaTeX的主要功能和使用方法。
\end{abstract}

\section{介绍}
\label{sec:introduction}

LaTeX是一个基于TeX的排版系统，特别适合生成高质量的技术和数学文档\cite{lamport1994latex}。与普通的文字处理软件相比，它使用标记命令来控制文档格式，能够更好地处理复杂的数学公式和专业排版。

\subsection{主要特点}
LaTeX具有以下主要特点：
\begin{itemize}
    \item Powerful mathematical formula typesetting capabilities
    \item Automatic generation of table of contents, cross-references, and bibliographies
    \item Highly customizable document styles
    \item Support for complex table and figure排版
    \item Suitable for collaborative editing of large documents
\end{itemize}

\section{数学公式}
\label{sec:math}

\subsection{行内公式}
行内公式如 $a^2 + b^2 = c^2$（勾股定理）和 $e^{i\pi} + 1 = 0$（欧拉恒等式）可以直接嵌入文本中。

\subsection{独立公式}
复杂的公式通常单独成行显示：
\begin{equation}
    \int\_{-\infty}^{\infty} e^{-x^2} dx = \sqrt{\pi}
    \label{eq:gaussian}
\end{equation}

\subsection{多行公式}
使用align环境进行多行公式的排版：
\begin{align}
    \sin(A+B) &= \sin A \cos B + \cos A \sin B \\
    \cos(A+B) &= \cos A \cos B-\sin A \sin B
\end{align}

\subsection{矩阵}
矩阵表示示例：
\[
\mathbf{M} =
\begin{pmatrix}
    m_{11} & m_{12} & \cdots & m_{1n} \\
    m_{21} & m_{22} & \cdots & m_{2n} \\
    \vdots & \vdots & \ddots & \vdots \\
    m_{m1} & m_{m2} & \cdots & m_{mn}
\end{pmatrix}
\]

\section{定理和定义}
\label{sec:theorems}
\begin{definition}
设$f(x)$在区间$[a,b]$上定义。如果对于任意$\epsilon > 0$，存在$\delta > 0$使得当$|x-x_0| < \delta$时，有$|f(x)-f(x_0)| < \epsilon$，则称$f(x)$在$x_0$处连续。
\end{definition}
\begin{theorem}
If a function $f(x)$ is continuous on the closed interval $[a,b]$, then $f(x)$ is uniformly continuous on $[a,b]$.
\end{theorem}
\begin{lemma}
如果数列$\{a_n\}$单调且有界，则$\{a_n\}$必定收敛。
\end{lemma}

\section{图形和表格}
\label{sec:figures}
\begin{figure}[htbp]
    \centering
\begin{subfigure}[b]{0.45\textwidth}
        \centering
        \includegraphics[width=\textwidth]{example-image-a}
        \caption{Sample image A}
        \label{fig:sub1}
\end{subfigure}
    \hfill
\begin{subfigure}[b]{0.45\textwidth}
        \centering
        \includegraphics[width=\textwidth]{example-image-b}
        \caption{Sample image B}
        \label{fig:sub2}
\end{subfigure}
    \caption{Example of subfigures}
    \label{fig:main}
\end{figure}

表\ref{tab:example}展示了一个简单的数据对比。
\begin{table}[h!]
    \centering
\begin{tabular}{|c|c|c|}
        \hline
        Method & Accuracy (\%) & Runtime (s) \\
        \hline
        A & 85.2 & 12.3 \\
        B & 90.5 & 18.7 \\
        C & 92.1 & 22.5 \\
        \hline
\end{tabular}
    \caption{Performance comparison of different methods}
    \label{tab:example}
\end{table}NK\_1071F892\_\_

表\ref{tab:example}显示了一个简单的数据比较。
\begin{table}[h!]
    \centering
\begin{tabular}{|c|c|c|}
        \hline
        Method & Accuracy (\%) & Runtime (s) \\
        \hline
        A & 85.2 & 12.3 \\
        B & 90.5 & 18.7 \\
        C & 92.1 & 22.5 \\
        \hline
\end{tabular}
    \caption{Performance comparison of different methods}
    \label{tab:example}
\end{table}

\section{算法}
\label{sec:algorithms}

以下算法\ref{alg:euclidean}描述了用于寻找最大公约数的欧几里得算法。
\begin{algorithm}
    \caption{欧几里得算法}
    \label{alg:euclidean}
\begin{algorithmic}[1]
        \Function{GCD}{$a, b$}
            \While{$b \neq 0$}
                \State $r \gets a \mod b$
                \State $a \gets b$
                \State $b \gets r$
            \EndWhile
            \Return $a$
        \EndFunction
\end{algorithmic}
\end{algorithm}

\section{代码示例}
\label{sec:code}

下面展示了一个计算阶乘的简单Python函数：
\begin{lstlisting}
def factorial(n):
    """Calculate the factorial of a non-negative integer"""
    if n < 0:
        raise ValueError("Factorial is not defined for negative numbers")
    result = 1
    for i in range(1, n+1):
        result *= i
    return result

# Example usage
print(factorial(5))  # Output: 120
\end{lstlisting}

\section{结论}
\label{sec:conclusion}

本文档涵盖了LaTeX排版的各种基本要素。通过掌握这些基本语法和环境，可以创建专业且格式良好的技术文档。有关更多高级功能，请参阅全面的LaTeX文档及相关包。

\bibliographystyle{plainnat}
\bibliography{references}
\end{document}